\section{From block counts to generating tuples}\label{sec:group}

Fix a prime $p$. Write $O_p(X)$ for the largest normal $p$-subgroup
of a finite group $X$, and put $[x,y]=xyx^{-1}y^{-1}$. A
$p$-element includes the identity; a $p'$-element has order prime to
$p$; and a $p$-singular element has order divisible by $p$.
The geometric sections will establish the following precise counting
statement. Its uniformity is part of the assertion.

\begin{theorem}[Uniform divisibility for generating tuples]
\label{thm:tuple-divisibility}
Let $J$ be a finite group with $O_p(J)=1$. For $s\geq0$, $n=p^s$,
and a multiset
$\mathcal M$ of $n$ conjugacy classes of $p$-singular elements of $J$,
let $N_{J,s}(\mathcal M)$ be the number, modulo simultaneous
$J$-conjugation, of ordered tuples
\[
 (x,y,u_1,\ldots,u_n)\in J^{n+2}
\]
satisfying
\begin{equation}\label{eq:tuple-relation}
 [x,y]u_1\cdots u_n=1,\qquad
 \langle x,y,u_1,\ldots,u_n\rangle=J,
 \qquad \{[u_1],\ldots,[u_n]\}=\mathcal M.
\end{equation}
For every integer $h\geq1$, there is an integer $s_0=s_0(p,J,h)$
such that
\[
 p^h\mid N_{J,s}(\mathcal M)
 \qquad(s\geq s_0)
\]
for every such multiset $\mathcal M$.
\end{theorem}

The entries $u_i$ remain ordered: prescribing their multiset of classes
does not quotient by permutations. Empty tuple sets have count zero,
so the statement includes $J=1$ and groups of order prime to $p$.
The proof of \Cref{thm:tuple-divisibility} is completed after
\Cref{prop:no-small-cover}. This section proves that it implies the
block equality, using \Cref{foundation:groups,foundation:algebra}
and proving the required counting reductions explicitly.

\subsection{Coefficient projection and the weight transform}

For each ambient finite group $X$, choose a splitting coefficient system
$(K,\mathcal O,k)$ as in \Cref{sec:introduction,sec:coefficients}, with
$X$ in place of $G$. Hold this system fixed throughout the calculations
with $X$ and its subquotients.
This system also splits all subquotients that occur below: a subgroup of a
quotient lifts to a subgroup by taking its preimage, and a
representation of a quotient inflates to that preimage. These
observations apply after any number of subgroup and quotient steps.
Ordinary irreducibles may therefore be realized over $K$. Their
identification with complex irreducibles, when needed, follows by
first splitting the group algebra over $\overline{\mathbb Q}$ and
then extending scalars to algebraically closed characteristic-zero
fields.

For a central idempotent $E\in kX$, define
\begin{align*}
 l(X,E)&=\#\{\text{simple $kX$-modules selected by $E$}\},\\
 k_{\mathrm{ord}}(X,E)&=\#\{\text{ordinary irreducibles selected by the lift of $E$}\},\\
 z(X,E)&=\#\{\chi\text{ selected by that lift}:\chi(1)_p=|X|_p\}.
\end{align*}
Here and throughout, modules and characters are counted up to
isomorphism. The class sums form a basis of $Z(\mathcal O X)$, and
their reductions form a basis of $Z(kX)$. Idempotent lifting over the
complete DVR thus gives a unique central idempotent lifting $E$.
Every integral central element acts on an ordinary irreducible by
a scalar in $\mathcal O$: the scalar lies in $K$ by splitting and
satisfies a monic integral equation by finiteness, and $\mathcal O$
is integrally closed. Reduction of this scalar defines a character
of $Z(kX)$. The Artinian local decomposition of that center shows
that it is precisely $\lambda_b$, where $b$ is the block containing
the ordinary irreducible.

For $H\leq X$, let $s_H$ denote coefficient restriction from
$Z(kX)$ to $Z(kH)$, or its integral analogue when appropriate.
For a $p$-subgroup $Q$, let $\operatorname{Br}_Q$ denote restriction
of coefficients from the $Q$-invariants $(kX)^Q$ to $kC_X(Q)$.

\begin{lemma}[Normalizer projection]\label{lem:normalizer-projection}
The map $\operatorname{Br}_Q:(kX)^Q\longrightarrow kC_X(Q)$ is a
unital algebra homomorphism. Suppose that
\[
 Q\leq H\leq N_X(Q),\qquad C_X(Q)\leq H.
\]
Writing $\pi_Q:kH\longrightarrow k[H/Q]$ for the quotient map, one
has, for $v\in Z(kX)$,
\begin{equation}\label{eq:normalizer-projection}
 \pi_Q(s_H(v))=\pi_Q(\operatorname{Br}_Q(v)),
 \qquad
 \lambda_b(s_H(v))=\lambda_b(\operatorname{Br}_Q(v))
\end{equation}
for every block $b$ of $H$. In particular $\lambda_b\circ s_H$ is
the central character of a unique block of $X$.
\end{lemma}

\begin{proof}
To compute the coefficient of $g\in C_X(Q)$ in a product of two
$Q$-invariant elements, let $Q$ conjugate the pairs $(a,b)$ with
$ab=g$. Coefficient products are constant on the orbits. Every
nontrivial orbit has size divisible by $p$, and the fixed pairs are
exactly those with $a,b\in C_X(Q)$. This proves multiplicativity;
the identity is preserved.

Each coset $hQ$ is stable under conjugation by $Q$, because
$h\in N_X(Q)$. Applying the same orbit cancellation to the sum
of coefficients over that coset proves the first identity in
\eqref{eq:normalizer-projection}. Also
$\operatorname{Br}_Q(v)\in Z(kH)$, since its support lies in
$C_X(Q)\leq H$ and it is invariant under $H$.

Every nonzero module for a finite $p$-group over $k$ has nonzero
invariants. For completeness, choose a central element $z$ of order
$p$; the nonzero kernel of $z-1$, whose $p$th power is zero, is a
module for the smaller group $Q/\langle z\rangle$, and induction
applies. If the module is simple for $H$, its $Q$-invariants are
$H$-stable, and hence the whole module. Thus every simple $kH$-module
factors through $H/Q$, proving the second identity. The composition
with $\operatorname{Br}_Q$ is a unital algebra character. Characters
of the finite commutative algebra $Z(kX)$ are its local residue
characters, one for each block, giving uniqueness.
\end{proof}

For $H=N_X(Q)$ put
\[
 \overline E_Q
   =\pi_Q(s_H(E))
   =\pi_Q(\operatorname{Br}_Q(E))\in Z(k[H/Q]).
\]
This is an idempotent. Its selected defect-zero quotient characters
are exactly the weights assigned to the blocks appearing in $E$.
Indeed, the lift of $\operatorname{Br}_Q(E)$ maps to the lift of
$\overline E_Q$ by uniqueness of central idempotent lifting. Its
action on the inflation of a quotient character is therefore one
exactly when the latter is selected by $\overline E_Q$. If $b$ is
the block of this inflation, the same test is
\[
 \lambda_b(\operatorname{Br}_Q(E))
   =\lambda_b(s_H(E))=1.
\]
By \Cref{lem:normalizer-projection}, this says that the induced block
is one of the blocks in $E$, with the coefficient-projection
definition in \eqref{eq:block-assignment}. Once $Q$ is fixed up to $X$-conjugacy,
there is no further identification of its quotient characters:
$N_X(Q)$ acts on $N_X(Q)/Q$ by inner automorphisms.

For a conjugation-invariant function $g$ of such pairs $(X,E)$,
define
\begin{equation}\label{eq:weight-transform}
 (Tg)(X,E)=
 \sum_{\substack{1<Q\leq X\text{ a $p$-subgroup}\\
                  \text{up to }X\text{-conjugacy}}}
 g\bigl(N_X(Q)/Q,\overline E_Q\bigr).
\end{equation}
Consequently the desired universal equality for block sums is
$l=(1+T)z$. The identity term includes precisely $Q=1$.
Every transition in $T$ strictly decreases the $p$-part of the group
order, so $T$ is nilpotent on each evaluated pair. Thus the formal
inverse $(1+T)^{-1}=\sum_{r\geq0}(-T)^r$ is evaluation-wise finite.

The next step makes the connection with subgroup-chain formulations of
the weight conjecture \cite{KnorrRobinson1989}. We derive the particular
normal-chain identity from this finite inverse, retaining its idempotent
and orbit data throughout.

\begin{definition}[Chains]\label{def:p-chains}
A chain in $X$ is a strictly increasing sequence of $p$-subgroups
\[
 \mathcal C=(1=Q_0<Q_1<\cdots<Q_r),
 \qquad |\mathcal C|=r.
\]
The length-zero chain $(1)$ is included. Set
\[
 H_{\mathcal C}=\bigcap_{i=0}^rN_X(Q_i),
 \qquad E_{\mathcal C}=\operatorname{Br}_{Q_r}(E).
\]
Call $\mathcal C$ \emph{normal} if every $Q_i$ is normal in $Q_r$.
\end{definition}

For every chain, normal or otherwise,
$C_X(Q_r)\leq H_{\mathcal C}\leq N_X(Q_r)$, so
$E_{\mathcal C}$ is an idempotent of $Z(kH_{\mathcal C})$.
For a normal chain one also has $Q_r\leq H_{\mathcal C}$.

\begin{lemma}[Finite chain inversion]\label{lem:chain-transform}
The identity $l=(1+T)z$ for all pairs $(X,E)$ is equivalent to
\begin{equation}\label{eq:chain-inversion}
 z(X,E)=
 \sum_{\substack{\mathcal C/X\\\mathcal C\ \mathrm{normal}}}
 (-1)^{|\mathcal C|}l(H_{\mathcal C},E_{\mathcal C}),
\end{equation}
where $\mathcal C/X$ denotes one representative of each
$X$-conjugacy orbit.
\end{lemma}

\begin{proof}
We identify the terms of $T^r l$. After $r$ steps they correspond
to normal chains of length $r$, with group
$H_{\mathcal C}/Q_r$ and idempotent
\[
 \pi_{Q_r}(s_{H_{\mathcal C}}(E))
   =\pi_{Q_r}(E_{\mathcal C}).
\]
For the induction, the preimage $Q_{r+1}$ of a nontrivial
$p$-subgroup of $H_{\mathcal C}/Q_r$ is a $p$-group strictly
containing $Q_r$. It lies in $H_{\mathcal C}$ and therefore
normalizes every preceding $Q_i$. Conversely every normal one-step
extension arises in this way. Its normalizer in the quotient has
preimage
$H_{\mathcal C}\cap N_X(Q_{r+1})=H_{\mathcal C'}$.
After fixing $\mathcal C$, conjugation is by its stabilizer
$H_{\mathcal C}$, which is exactly quotient conjugation at the next
step. This proves the orbit assertion without an extra multiplicity.

To check the idempotent at the next step, write
$E=\sum_{x\in X}e_xx$ and put $H'=H_{\mathcal C'}$.
The new normalizer in $H_{\mathcal C}/Q_r$ is $H'/Q_r$.
Restriction to this normalizer retains exactly the coefficient sums
over the $Q_r$-cosets contained in $H'$. Choose a set $A$ of
representatives for the cosets of $Q_r$ in $Q_{r+1}$. For $h\in H'$,
the subsequent quotient has coefficient
\[
 \sum_{a\in A}\ \sum_{q\in Q_r}e_{haq}
     =\sum_{t\in Q_{r+1}}e_{ht}
\]
at $hQ_{r+1}$: the sets $aQ_r$, $a\in A$, partition $Q_{r+1}$.
These are precisely the coefficients of
$\pi_{Q_{r+1}}(s_{H'}(E))$.
By \Cref{lem:normalizer-projection}, this idempotent is
$\pi_{Q_{r+1}}(E_{\mathcal C'})$, as required.
Every simple $kH_{\mathcal C}$-module factors through
$H_{\mathcal C}/Q_r$ and has the indicated idempotent action.
Hence the contribution to $T^rl$ is the corresponding term in
\eqref{eq:chain-inversion}. The finite inverse of $1+T$ now gives
the equivalence.
\end{proof}

\subsection{Cocenters and subsections}

By \Cref{lem:chain-transform}, the remaining algebraic objective is
\eqref{eq:chain-inversion}. We next express ordinary character counts
as sums of simple-module counts in centralizers. Chain cancellation
will then replace the alternating sum of $l$ by an alternating sum of
$k_{\mathrm{ord}}$, which can be computed by central traces.

For an algebra $A$, write $V(A)=A/[A,A]$, where $[A,A]$ is the
linear span of commutators, and write $\overline a$ for the image of
$a$. A central element acts on $V(A)$ by multiplication.

Powers modulo commutators underlie K\"ulshammer's generalized Reynolds
ideals; see \cite[Section~2]{HHKM2005}. We give the finite-dimensional
argument needed to extract the two block counts.

\begin{lemma}[Cocenter ranks]\label{lem:cocenter-ranks}
For a finite-dimensional $k$-algebra $A$, the map
\[
 F_A:V(A)\longrightarrow V(A),\qquad
 \overline a\longmapsto\overline{a^p}
\]
is well-defined and Frobenius semilinear. Its stable image has
dimension equal to the number of simple $A$-modules, and its
summand selected by a central idempotent has dimension equal to
the corresponding simple-module count.
For $A=kX$, this stable image is the span of the $p'$-classes.
The ranks of multiplication by $E$ on that span and on all of
$V(kX)$ are, respectively, $l(X,E)$ and $k_{\mathrm{ord}}(X,E)$.
\end{lemma}

\begin{proof}
In the expansion of $(a+b)^p$, cyclic rotations of any mixed word
give the same cocenter class. Their orbits have size $p$, leaving
only $a^p$ and $b^p$. Thus the power operation into the cocenter
is additive and is Frobenius semilinear. Moreover
$\overline{(ab)^p}=\overline{(ba)^p}$, by one cyclic rotation;
additivity shows that this operation kills the linear commutator
space. It therefore descends to $F_A$.

Let $R$ be the Jacobson radical of $A$. The kernel of
$V(A)\longrightarrow V(A/R)$ is represented by elements of $R$:
lift any commutator expression downstairs and subtract it upstairs.
Since $R$ is nilpotent, a sufficiently high iterate of $F_A$ kills
this kernel. On the split semisimple algebra $A/R$, the cocenter
has one matrix-trace coordinate per simple factor, and $F_A$
raises each coordinate to the $p$th power. It is bijective because
$k$ is perfect. For example, the trace identity here follows by
triangularizing a matrix over $k$ and taking the $p$th powers of
its diagonal entries.

The descending images of $F_A$ stabilize, since they are
$k$-subspaces of a finite-dimensional space. Denote the stable image
by $S$ and the quotient map by $q:V(A)\to V(A/R)$.
Choose $m$ large enough that $S=\operatorname{im}F_A^m$ and
$F_A^m$ kills $\ker q$. On $S$, the map $F_A$ is surjective
and hence bijective, because it is semilinear for the automorphism
$a\mapsto a^p$ of the perfect field $k$. If $x\in S\cap\ker q$,
then $F_A^m(x)=0$; injectivity on $S$ gives $x=0$.
To prove surjectivity of $q|_S$, take $y\in V(A/R)$. The
Frobenius map $F_{A/R}$ is bijective, so choose $y'$ with
$F_{A/R}^m(y')=y$ and lift $y'$ to $x\in V(A)$. Then
$F_A^m(x)\in S$ and $q(F_A^m(x))=y$. Thus $q|_S$ is an
isomorphism. Finally
$(Ea)^p=Ea^p$ for a central idempotent $E$, so all these statements
respect its summand.

For a group algebra the cocenter is free on conjugacy classes,
represented by individual group elements: the commutator
relations identify exactly conjugate elements. Sufficiently high
$p$-powers land on $p'$-elements, and taking $p$th powers permutes
the $p'$-classes. This identifies the stable image. The integral
cocenter $V(\mathcal O X)$ is free on the same class basis. The
lift of $E$ projects it onto a free direct summand, whose rank is
unchanged over $k$ and $K$. Split semisimplicity over $K$ identifies
this rank with $k_{\mathrm{ord}}(X,E)$.
\end{proof}

Every group element has a unique commuting factorization into its
$p$-part and $p'$-part, both powers of that element. A
\emph{section} is the set of elements whose $p$-parts are conjugate
to a fixed $p$-element.

The next identity is Brauer's classical subsection count
\cite{Brauer1959}. The cocenter argument also proves the support
property needed for the defect-zero detector, without dividing by
class sizes in characteristic $p$.

\begin{lemma}[Subsection formula and support]\label{lem:subsections}
For every central idempotent $E\in kX$,
\begin{equation}\label{eq:subsections}
 k_{\mathrm{ord}}(X,E)=
 \sum_{\substack{[u]\text{ an }X\text{-class}\\u\text{ a }p\text{-element}}}
 l\bigl(C_X(u),\operatorname{Br}_{\langle u\rangle}(E)\bigr).
\end{equation}
Multiplication by $E$ preserves the section decomposition of the
cocenter and the section decomposition of the center into spans
of class sums.
\end{lemma}

\begin{proof}
Fix a $p$-element $u$ and set $C=C_X(u)$. The map
$\overline t\mapsto\overline{ut}$ identifies the $p'$-class span
of $V(kC)$ with the section span belonging to $u$ in $V(kX)$.
It is a bijection of class bases: conjugacy of $ut$ and $ut'$
forces a conjugating element to centralize their common $p$-part
$u$, and the converse is immediate.

For $t\in C$ a $p'$-element, conjugation by $\langle u\rangle$
in the expansion of $Eut$ gives
\[
 \overline{Eut}
   =\overline{\operatorname{Br}_{\langle u\rangle}(E)ut}.
\]
Indeed, the coefficients are constant on these orbits and the
cocenter classes of their products with $ut$ agree; only fixed
orbits survive. The restricted idempotent is central in $kC$ and
preserves its $p'$-class span by \Cref{lem:cocenter-ranks}.
Multiplication by $u$ is central in $kC$, so the displayed
identification intertwines the two idempotent actions. Their ranks
are thus the terms in \eqref{eq:subsections}. Summing the ranks
over the sections proves that identity.

For the center, use the pairing
\[
 Z(kX)\times V(kX)\longrightarrow k,
 \qquad (a,\overline b)\longmapsto\tau(ab),
\]
where $\tau$ takes the coefficient of $1$. It is nondegenerate:
a class sum pairs with an individual representative of its inverse
class with coefficient one. Multiplication by $E$ is self-adjoint
for this pairing. Sections on the center pair with the inverse
sections of the cocenter, so preservation there proves preservation
here. In particular no class-size division is involved.
\end{proof}

\subsection{Two cancellations on chain orbits}

\begin{lemma}[Removal of nonnormal chains]\label{lem:nonnormal-cancellation}
In an alternating sum over chain orbits, every contribution that
depends only on $H_{\mathcal C}$ and $Q_r$ cancels on the nonnormal
chains. Thus such a sum may be taken over all chains or only normal
chains.
\end{lemma}

\begin{proof}
For a nonnormal chain, let $i$ be the largest index for which
$A=Q_i$ is not normal in $P=Q_r$, and let $R$ be the normal closure
of $A$ in $P$. Then $A<R\leq Q_{i+1}$: every later member is
normal in $P$. Also $R<P$. To see the latter, put the proper
subgroup $A$ in a maximal subgroup of $P$, which is normal because
$P$ is a finite $p$-group; it contains $R$.

Insert $R$ immediately after $A$ if it is absent, and remove it
if it is present. The last group $P$ and the last nonnormal member
$A$ are unchanged, so this operation is an involution. Every
element normalizing $A$ and $P$ normalizes their normal closure
$R$; hence insertion or removal leaves $H_{\mathcal C}$ unchanged.
The operation commutes with $X$-conjugation and changes the length
by one. It induces a fixed-point-free sign reversal on the orbits,
which proves cancellation.
\end{proof}

\begin{lemma}[Alternating subsection cancellation]
\label{lem:alternating-subsections}
One has
\begin{equation}\label{eq:alternating-subsections}
 \sum_{\substack{\mathcal C/X\\\mathcal C\ \mathrm{normal}}}
 (-1)^{|\mathcal C|}l(H_{\mathcal C},E_{\mathcal C})
 =
 \sum_{\substack{\mathcal C/X\\\mathcal C\ \mathrm{normal}}}
 (-1)^{|\mathcal C|}k_{\mathrm{ord}}(H_{\mathcal C},E_{\mathcal C}).
\end{equation}
\end{lemma}

\begin{proof}
By \Cref{lem:nonnormal-cancellation}, both sums may be taken over
all chains. Expand the ordinary count by \Cref{lem:subsections}.
The $u=1$ terms give the simple-module sum. The other terms are
indexed by $X$-orbits of pairs $(\mathcal C,u)$, with $u\ne1$ a
$p$-element in $H_{\mathcal C}$. This orbit description follows
because the stabilizer of a fixed chain is $H_{\mathcal C}$.
The contribution of such a pair is
\begin{equation}\label{eq:pair-contribution}
 l\bigl(H_{\mathcal C}\cap C_X(u),
         \operatorname{Br}_{Q_r\langle u\rangle}(E)\bigr).
\end{equation}
Here $u$ normalizes $Q_r$, so $Q_r\langle u\rangle$ is a
$p$-group; successive coefficient restrictions have support
$C_X(Q_r)\cap C_X(u)=C_X(Q_r\langle u\rangle)$, proving the
idempotent in this expression.

Let $i$ be the largest index, allowing $0$ and $r$, for which
$u\notin Q_i$. Such an index exists since $Q_0=1$. Toggle
$R=Q_i\langle u\rangle$ immediately after $Q_i$: insert it
if absent, and remove it if present. It is a $p$-group, since
$u$ normalizes $Q_i$, and it is contained in every later member,
all of which contain $u$. The last member not containing $u$
therefore remains $Q_i$; this proves involutivity, including an
insertion after the terminal member and its inverse deletion.

Every element centralizing $u$ and normalizing $Q_i$ normalizes
$R$. Thus $H_{\mathcal C}\cap C_X(u)$ is unchanged. The product
of the terminal member with $\langle u\rangle$ is also unchanged,
even when the terminal member is inserted or deleted. Hence
\eqref{eq:pair-contribution} is preserved. The operation is
$X$-equivariant and reverses the sign, so all $u\ne1$ terms cancel.
\end{proof}

\subsection{Detecting defect zero by central elements}

By \Cref{lem:alternating-subsections}, it remains to prove that the
alternating ordinary-character sum in
\eqref{eq:alternating-subsections} equals $z(X,E)$. Their difference
retains the signed contributions of all characters selected by
$E_{\mathcal C}$ on positive-length normal chains, but omits the
selected defect-zero characters at the chain $(1)$. We will realize
this difference as a trace limit by constructing a central element
with the corresponding zero-or-one central-character values.

Put
\[
 c_X=\sum_{u\in X\,\text{a }p\text{-element}}u,
 \qquad d_X=\sum_{x,y\in X}[x,y],
 \qquad I_X(a)=\sum_{x\in X}xax^{-1}.
\]
These elements and the image of $I_X$ are central. Call a block
\emph{semisimple} if its algebra is semisimple; since a block has no
nontrivial central idempotent, such a block is one full matrix
algebra over $k$.

The scalar of $d_X$ will identify these semisimple blocks with the
ordinary defect-zero characters. The complement $1-c_X$ will then remove
exactly those blocks from the central trace, using the support
information supplied by $I_X$.

The operator $I_X$ is the group-algebra Higman trace
\cite[Section~2]{Higman1955}; see also
\cite[Sections~2.4--2.5]{LorenzTokoly2011}. The use of $c_X$ to
isolate defect zero goes back to Tsushima \cite[Theorem~1]{Tsushima1971}.
The following proof records the exact radical-annihilation and
central-character consequences used in our trace limit.

\begin{lemma}[Averaging and defect zero]\label{lem:defect-zero-detection}
The image of $I_X$, and also $d_X$, annihilates the Jacobson radical
of $kX$. On a semisimple block, $d_X$ has nonzero central scalar;
on every other block its central character is zero. Each semisimple
block contains exactly one ordinary irreducible, which has defect
zero, and every ordinary defect-zero irreducible occurs in this
way. Moreover:
\begin{enumerate}[label=\textup{(\roman*)}]
 \item $\operatorname{im}I_X$ is the span of class sums of elements
 whose centralizers have order prime to $p$;
 \item each semisimple block idempotent belongs to
 $\operatorname{im}I_X$ and has zero Brauer image at every nontrivial
 $p$-subgroup;
 \item $\lambda_b(c_X)=0$ on every nonsemisimple block $b$.
\end{enumerate}
\end{lemma}

\begin{proof}
\textit{Radical annihilation.}
The form $(a,b)\mapsto\tau(ab)$ on $kX$ is symmetric and
nondegenerate. For dual bases $v_i,v_i^*$, the tensor
$\Omega=\sum_i v_i\otimes v_i^*$ is independent of the chosen
bases. Apply the linear map $u\otimes v\mapsto uav$ to this
tensor. In the group basis, whose dual basis consists of the inverses,
the result is $I_X(a)$. Similarly, apply
$u\otimes v\otimes w\otimes z\mapsto uwvz$ to
$\Omega\otimes\Omega$. In the group basis the result is
$\sum_{x,y\in X}xyx^{-1}y^{-1}=d_X$. Thus in any dual bases
\begin{equation}\label{eq:averaging-tensor}
 I_X(a)=\sum_i v_iav_i^*,\qquad
 d_X=\sum_{i,j}v_iv_jv_i^*v_j^*
     =\sum_j I_X(v_j)v_j^*.
\end{equation}
For $b$ in the radical $R$, the trace identity
\[
 \tau(I_X(a)b)=\operatorname{Tr}(L_bR_a)
\]
follows by writing the trace in these dual bases; $L$ and $R$
denote left and right multiplication. Since $L_b$ is nilpotent
and commutes with $R_a$, this trace is zero. Replace $b$ by $bc$
for arbitrary $c\in kX$, and use nondegeneracy to obtain
$I_X(a)b=0$. The final expression in
\eqref{eq:averaging-tensor} then shows that $d_XR=0$ as well,
because $v_j^*R\subseteq R$.

\textit{Semisimple blocks and ordinary defect zero.}
Distinct block ideals are orthogonal for the form. On a
semisimple block $M_d(k)$ its restriction is $t\operatorname{Tr}$
for some $t\in k^\times$. Indeed a symmetric functional kills
commutators, whose quotient in a matrix algebra is the trace line,
and nondegeneracy forces $t\ne0$. Matrix units $e_{ij}$ have duals
$t^{-1}e_{ji}$. Since
$\sum_{i,j}e_{ij}ae_{ji}=\operatorname{Tr}(a)1$, the first
contraction in \eqref{eq:averaging-tensor} gives
$I_X(a)=t^{-1}\operatorname{Tr}(a)1$. In particular
$I_X(e_{ij})=t^{-1}\delta_{ij}1$, so its last expression gives
\[
 d_X=\sum_{i,j}I_X(e_{ij})t^{-1}e_{ji}
     =t^{-2}\sum_i e_{ii}=t^{-2}1.
\]
We have therefore proved
\begin{equation}\label{eq:simple-block-scalars}
 I_X(a)=t^{-1}\operatorname{Tr}(a)1,
 \qquad d_X=t^{-2}1
 \quad\text{on this block}.
\end{equation}
In particular its idempotent lies in the image of $I_X$; one can
take $a=te_{11}$. These formulas do not require $d$ to be nonzero
in $k$. On a nonsemisimple block, a central element with nonzero
residue scalar is a unit in its local center. Such a unit cannot
annihilate the nonzero radical. This proves the assertion about
the central character of $d_X$.

Now compute the same commutator sum in characteristic zero. If
$\chi$ has degree $d$, Schur's lemma gives
\[
 \sum_{x\in X}\rho(xyx^{-1})
      =\frac{|X|\chi(y)}d\,1.
\]
Multiply by $\rho(y^{-1})$, sum over $y$, and take traces.
Ordinary character orthogonality gives the scalar
\begin{equation}\label{eq:ordinary-commutator-scalar}
 \lambda_\chi(d_X)=\frac{|X|^2}{\chi(1)^2}.
\end{equation}
It is integral, by the integrality of ordinary central scalars.
Its reduction is nonzero exactly when
$v_p(\chi(1))=v_p(|X|)$, namely in defect zero. The modular
calculation therefore puts exactly these characters in
semisimple blocks. Such a block has cocenter dimension one, so
\Cref{lem:cocenter-ranks} shows that its lifted block selects
exactly one ordinary irreducible. This also proves existence of
that character.

\textit{Support and the $p$-element sum.}
For a group element $g$ one has
\[
 I_X(g)=|C_X(g)|\sum_{a\in[g]}a.
\]
This proves assertion (i). The support of a semisimple block
idempotent, already in this image, therefore contains no element
centralized by a nontrivial $p$-subgroup, proving (ii).

Finally fix a nontrivial $p$-subgroup $P$. Choose
$z\in Z(P)$ of order $p$. It is central also in $C_X(P)$, and
multiplication by $z$ preserves the set of $p$-elements of that
centralizer. With $q=\sum_{j=0}^{p-1}z^j$, the sum of those
$p$-elements factors as $q$ times a sum of orbit representatives.
As $q$ is central and $q^2=pq=0$, this proves
\[
 \operatorname{Br}_P(c_X)^2=0,
 \qquad \operatorname{Br}_P(c_X^2)=0.
\]
If $p$ divides $|C_X(g)|$, a nontrivial $p$-subgroup
$P\leq C_X(g)$ makes the coefficient of $g$ in $c_X^2$ vanish
by this identity. Assertion (i) gives $c_X^2\in\operatorname{im}I_X$.
On a nonsemisimple block its central character must therefore
vanish, by the same radical-annihilator argument. Since
$\lambda_b(c_X)^2=\lambda_b(c_X^2)=0$, assertion (iii) follows.
\end{proof}

\begin{lemma}[An element with zero-or-one residues]
\label{lem:detecting-element}
Let $v=E(1-c_X)$. It is central and supported on $p$-singular
elements. For a normal chain $\mathcal C$ of positive length, its
restriction to $H_{\mathcal C}$ satisfies
\[
 \lambda_b(s_{H_{\mathcal C}}(v))=\lambda_b(E_{\mathcal C})
 \in\{0,1\}
\]
for every block $b$ of $H_{\mathcal C}$. For the length-zero
chain, the scalar is one on the nonsemisimple blocks selected by
$E$ and zero on every other block.
\end{lemma}

\begin{proof}
The element $1-c_X$ is the negative sum of the nonidentity
$p$-elements, hence lies entirely in $p$-singular sections.
Section preservation from \Cref{lem:subsections} proves the same
support assertion for $v$. Let $B_0$ be a semisimple block
idempotent. Since $v$ is central and $B_0kX$ is a full matrix
algebra, $B_0v$ lies in its one-dimensional center $kB_0$.
On the other hand, section preservation makes the support of $B_0v$
$p$-singular. By \Cref{lem:defect-zero-detection}, the support
of $B_0$ contains only elements whose centralizers have order prime
to $p$, and therefore no $p$-singular elements. The two support
conditions force the scalar multiple $B_0v$ to be zero.
On a nonsemisimple block,
\Cref{lem:defect-zero-detection} gives
$\lambda_b(v)=\lambda_b(E)$, proving the length-zero assertion.

For positive length, put $Q=Q_r\ne1$ and $H=H_{\mathcal C}$.
By \Cref{lem:normalizer-projection}, $\lambda_b\circ s_H$ is a
central character of $X$ and factors through
$\operatorname{Br}_Q$. It cannot be the character of a
semisimple block $B_0$: that would give value one on $B_0$,
whereas $\operatorname{Br}_Q(B_0)=0$ by
\Cref{lem:defect-zero-detection}. On its nonsemisimple block,
$\lambda(v)=\lambda(E)$. Applying
\Cref{lem:normalizer-projection} to $E$ gives the asserted value
$\lambda_b(E_{\mathcal C})$.
\end{proof}

\subsection{Central traces and the final cancellation}

Choose a class-sum lift
\[
 \widehat v=\sum_{a\in X}\alpha_a a\in Z(\mathcal O X)
\]
of the detecting element $v=E(1-c_X)$ from
\Cref{lem:detecting-element}, with $\alpha_a=0$ whenever $a$ is
not $p$-singular.
For a positive integer $n$, define
\begin{equation}\label{eq:chain-trace}
 \Phi_n=
 \sum_{\mathcal C/X}(-1)^{|\mathcal C|}
 \operatorname{Tr}\left(
    s_{H_{\mathcal C}}(\widehat v)^n
       \mid V(KH_{\mathcal C})\right),
\end{equation}
where the operator is central multiplication and all chains are
included.

\begin{lemma}[The first trace limit]\label{lem:trace-limit}
In the valuation topology of $K$,
\begin{equation}\label{eq:trace-limit}
 \lim_{s\to\infty}\Phi_{p^s}
 =
 \sum_{\substack{\mathcal C/X\\\mathcal C\ \mathrm{normal}}}
 (-1)^{|\mathcal C|}k_{\mathrm{ord}}(H_{\mathcal C},E_{\mathcal C})
 -z(X,E).
\end{equation}
\end{lemma}

\begin{proof}
The nonnormal chains cancel by
\Cref{lem:nonnormal-cancellation}, since their trace contribution
depends only on the stabilizer. For a remaining chain, the
cocenter of $KH_{\mathcal C}$ has one coordinate per ordinary
irreducible. The unpowered operator has on these coordinates
integral central eigenvalues reducing to the scalars in
\Cref{lem:detecting-element}.

If an eigenvalue reduces to zero, its $p^s$th powers tend to zero.
If it is $1+a$ with $\operatorname{val}(a)>0$, the binomial
formula gives
\[
 \operatorname{val}((1+a)^p-1)
 \geq
 \min\{\operatorname{val}(p)+\operatorname{val}(a),
                    p\operatorname{val}(a)\}.
\]
After normalizing the discrete valuation to integer values, each
iteration increases the positive valuation by at least one, unless
the difference has already become zero. These powers tend to one.
This argument applies also to $p=2$ and requires no finiteness of
the residue field.

For positive-length chains the limiting trace therefore counts
all ordinary irreducibles selected by $E_{\mathcal C}$. At length
zero it omits the semisimple blocks, each of which contributes
exactly one defect-zero character by
\Cref{lem:defect-zero-detection}. There are finitely many chains
and irreducibles, so limits may be summed, proving
\eqref{eq:trace-limit}.
\end{proof}

The next identity follows from Frobenius's commutator-counting formula
and is the weighted genus-one case of the surface formula in
\cite[Theorem~2.1 and Equation~(3.1)]{Klug2025}. Its single
denominator is important when chain orbits are replaced by actual
chains, so we verify the normalization directly.

\begin{lemma}[Trace normalization]\label{lem:tuple-trace}
For $H\leq X$ and $w=s_H(\widehat v)$,
\begin{equation}\label{eq:tuple-trace}
 \operatorname{Tr}(w^n\mid V(KH))
 =\frac1{|H|}
   \sum_{\substack{x,y,u_1,\ldots,u_n\in H\\
                   [x,y]u_1\cdots u_n=1}}
       \prod_{i=1}^n\alpha_{u_i}.
\end{equation}
\end{lemma}

\begin{proof}
The sum on the right before division is
$\tau(d_Hw^n)$. The regular trace of any element of $KH$ is
$|H|$ times its identity coefficient. On the other hand,
\eqref{eq:ordinary-commutator-scalar} and the regular decomposition
give
\begin{align*}
 |H|\tau(d_Hw^n)
  &=\sum_{\chi\in\operatorname{Irr}(H)}
       \chi(1)^2\frac{|H|^2}{\chi(1)^2}\lambda_\chi(w)^n\\
  &=|H|^2\operatorname{Tr}(w^n\mid V(KH)).
\end{align*}
Dividing proves exactly the single denominator in
\eqref{eq:tuple-trace}.
\end{proof}

For a subgroup $J\leq X$, define an integer on actual chains,
without quotienting by conjugation,
\begin{equation}\label{eq:chain-weight}
 a_X(J)=
 \sum_{\substack{\mathcal C\text{ a chain in }X\\
                  J\leq H_{\mathcal C}}}(-1)^{|\mathcal C|}.
\end{equation}

\begin{lemma}[Cancellation by a normal $p$-subgroup]
\label{lem:pcore-cancellation}
If $O_p(J)\ne1$, then $a_X(J)=0$.
\end{lemma}

\begin{proof}
Put $P=O_p(J)$. In any chain normalized by $J$, each $Q_i$ is
normalized by $P\leq J$, so $Q_iP$ is a $p$-subgroup. Both
factors are normalized by $J$, making their product
$J$-normalized. Let $i$ be the last index such that $P\nleq Q_i$;
it exists since $Q_0=1$ and $P\ne1$. Toggle $Q_iP$ immediately
after $Q_i$. It strictly contains $Q_i$ and is contained in all
later members. Insertion and deletion preserve the last member
not containing $P$, so this is an involution. It is defined also
for the length-zero chain and for terminal insertion or deletion.
It preserves the property of being normalized by $J$ and reverses
the sign. The sum of actual chains in \eqref{eq:chain-weight}
therefore vanishes.
\end{proof}

\begin{proposition}[Reduction to uniform tuple divisibility]
\label{prop:group-reduction}
If \Cref{thm:tuple-divisibility} holds for every finite group with
trivial $p$-core, then for every finite group $X$ and central
idempotent $E\in kX$ one has $l(X,E)=((1+T)z)(X,E)$.
In particular the block equality in \Cref{thm:main} follows, with
the stated block assignment and with $Q=1$ included.
\end{proposition}

\begin{proof}
The orbit of a chain $\mathcal C$ has size
$|X|/|H_{\mathcal C}|$. Consequently
\Cref{lem:tuple-trace} changes the orbit sum
\eqref{eq:chain-trace} into the following sum over actual tuples
and actual chains:
\begin{equation}\label{eq:trace-by-generated-group}
 \Phi_n=\frac1{|X|}
 \sum_{\substack{x,y,u_1,\ldots,u_n\in X\\
                 [x,y]u_1\cdots u_n=1}}
 a_X\bigl(\langle x,y,u_1,\ldots,u_n\rangle\bigr)
 \prod_{i=1}^n\alpha_{u_i}.
\end{equation}
To verify the factor, each orbit contribution has prefactor
$1/|H_{\mathcal C}|$; replacing its representative by all
$|X|/|H_{\mathcal C}|$ conjugates gives the common prefactor
$1/|X|$. For a fixed tuple the admissible chains are exactly those
normalized by its generated subgroup. This is precisely
$a_X(J)$.

By \Cref{lem:pcore-cancellation}, only generated subgroups
$J$ with $O_p(J)=1$ contribute. By the support of $\widehat v$,
only tuples with every $u_i$ $p$-singular contribute. Fix such a
subgroup $J$ and group its generating tuples according to the
multiset $\mathcal M$ of $J$-conjugacy classes of the $u_i$.
The product $\prod_i\alpha_{u_i}$ is constant on this multiset:
the coefficients are constant on $X$-classes, hence on $J$-classes,
and their product is unchanged by reordering. Denote this integral
product by $\alpha(\mathcal M)$. A generating tuple has simultaneous
conjugation stabilizer $Z(J)$, since an element centralizing all
its entries centralizes the group they generate. Thus its number
of actual tuples is $|J|/|Z(J)|$ times its orbit count. For
$n=p^s$, the numerator in \eqref{eq:trace-by-generated-group}
is therefore
\begin{equation}\label{eq:uniform-numerator}
 \sum_{\substack{J\leq X\\O_p(J)=1}}
 a_X(J)\frac{|J|}{|Z(J)|}
 \sum_{\mathcal M}
   \alpha(\mathcal M)N_{J,s}(\mathcal M).
\end{equation}

Fix any $h\geq1$. By \Cref{thm:tuple-divisibility}, each summand
of the inner sum belongs to $p^h\mathcal O$ for all sufficiently
large $s$, with a threshold independent of $\mathcal M$. Although
the number of multisets grows with $s$, the sum is finite for each
$s$ and remains in that same ideal. There are only finitely many
subgroups $J\leq X$, so take the largest of their thresholds.
The integer factors in \eqref{eq:uniform-numerator} preserve
divisibility. We conclude that its numerator tends to zero in the
valuation topology. The fixed denominator $|X|$ does not change
this conclusion, so $\Phi_{p^s}\longrightarrow0$.

Combine this with \Cref{lem:trace-limit} and
\Cref{lem:alternating-subsections}. The resulting equality is
\eqref{eq:chain-inversion}; its integer terms agree in the
characteristic-zero field $K$, and hence agree as integers.
The same reasoning applies to every pair and to the splitting
subquotients arising in the transform. \Cref{lem:chain-transform}
now gives $l=(1+T)z$. The weight interpretation of
\eqref{eq:weight-transform} proves the asserted block equality.
\end{proof}
